\chapter{Antecedentes}
\label{ch:antecedentes}

% Fuente: apartado «Antecedentes» de la propuesta aprobada.
% TODO: desarrollar cada sección con fuentes verificadas (ver TODO en referencias.bib).

El auge de los dispositivos \textit{wearables} (Garmin, Polar, Suunto) y de plataformas como Strava o TrainingPeaks ha impulsado la generación masiva de datos de rendimiento deportivo. Sin embargo, la mayoría de estos sistemas presentan interfaces cerradas o de uso limitado para análisis personalizados. El desarrollo de \textit{dashboards} propios mediante herramientas como Power BI, Python (Dash, Streamlit) o React permite crear soluciones adaptadas a las necesidades del usuario final, integrando múltiples fuentes y técnicas de análisis predictivo.

Este trabajo se enmarca dentro de las metodologías de gestión basada en datos (\textit{Data-Driven Management}) aplicadas al contexto deportivo \cite{davenport2010analytics, marr2016big}.

\section{Dispositivos y plataformas de registro deportivo}
\label{sec:plataformas}
% TODO: relojes GPS y sensores (FC, potencia, cadencia); plataformas (Strava, Garmin Connect,
% TrainingPeaks); qué análisis ofrecen y sus limitaciones (interfaces cerradas, exportación).

\section{Métricas de rendimiento en carrera}
\label{sec:metricas}
% TODO: ritmo y velocidad, \gls{fc} y zonas, cadencia, potencia.
% TODO: carga de entrenamiento (\gls{trimp}), relación carga aguda/crónica (\gls{acwr}),
% eficiencia aeróbica, predicción de marcas (fórmula de Riegel). Cada una con su fuente.

\section{Técnicas de análisis y predicción}
\label{sec:tecnicas}
% TODO: regresión lineal y series temporales \cite{montgomery2018applied}; aprendizaje
% automático aplicado \cite{geron2022hands}; validación temporal de modelos.

\section{Herramientas para el desarrollo de \textit{dashboards}}
\label{sec:herramientas_dashboard}
% TODO: comparativa Power BI, Streamlit, Dash, React (criterios: control de la interfaz,
% integración de modelos propios, versionado, coste). Enlaza con la decisión del ADR 0001.

\section{Conclusiones del estudio de antecedentes}
\label{sec:conclusiones_antecedentes}
% TODO: qué hueco cubre este trabajo frente a las soluciones existentes.
